%& -shell-escape
% !TeX root = thesis-text.tex
% !TeX encoding = UTF-8
% !TeX program = XeLaTeX

\documentclass{.local/lib/classes/ndvr-super-article-standalone}

\title[NDVR thesis-abstract]{%
    Проблема поиска нечетких дубликатов видео%
}

\begin{document}

    \section{Мониторинг эфира}

    Отслеживать видео по его содержимому 
    требуется не только для материалов, 
    размещенных в различных библиотеках видео.
    Аналогичная проблема стоит перед телеканалами.

    Особенно остро стоит проблема поиска 
    и замены эфирной рекламы \cite{Huang:2010b,Tan:2009,Yan:2008}.
    Для обеспечения максимального продающего эффекта
    реклама должна показываться с~правильной частотой,
    в~течение ожидаемого периода времени.

    Кроме того, с середины 2010-х годов 
    многие крупные каналы начали транслировать 
    свои визуальные потоки через Интернет. 
    Это сильно усложнило задачу мониторинга эфира. 

    Контроль локальных эфиров государственных каналов
    чаще всего осуществляется вручную.
    Для интернет-вещания какой-либо 
    контроль обычно и~вовсе отсутствует.
    И этим активно пользуются сайты,
    которые паразитируют на телевизионных эфирах.

    При~таких обстоятельствах для~мониторинга потока видео могут
    быть использованы методы обнаружения нечетких дубликатов.
    Когда искомый материал появится в~эфире, 
    методы обнаружения дубликатов смогут обнаружить 
    его возникновение и~сделать соответствующую запись.
    Таким образом можно будет проконтролировать количество показов,
    время показа, частоту и~другие параметры.

    % На~практике для~интернет-видео проблема решается гораздо проще.
    % Плеер, показывающий целевое видео, получает рекламу отдельно
    % и~отдельно ее проигрывает. Само целевое видео при~этом останавливается.
    % По~определенным событиям плеер отсылает некоторые запросы:
    % клик, начало показа, четверть показа, половина показа и~т.~д.
    % Эти запросы могут отсылаться в~том числе и~на~ресурс компании-заказчика.
    % Таким образом, поиск НДВ
    % не~является необходимым в~данном случае.
    % Нужно отметить, что~такой подход не~будет работать с~онлайн-трансляциями
    % телевизионных передач и~некоторыми другими видами потокового вещания.


\end{document}